\documentclass[12pt,a4paper]{article}

% Encoding, language, typography
\usepackage[T1]{fontenc}
\usepackage[utf8]{inputenc}
\usepackage[english]{babel}
\usepackage{lmodern}
\usepackage{microtype}

% Page layout
\usepackage[margin=2cm,top=3cm,bottom=3cm]{geometry}
\usepackage{parskip}
\usepackage{setspace}
\onehalfspacing

% Mathematics
\usepackage{amsmath}
\usepackage{amssymb}
\usepackage{amsfonts}
\usepackage{mathtools}
\usepackage{bm}

% Units and numbers
\usepackage{siunitx}

\sisetup{
    per-mode=symbol,
    separate-uncertainty=true,
    detect-all=true
}

% Quantum notation and circuits
\usepackage{braket}
\usepackage{quantikz}

% Figures and tables
\usepackage{graphicx}
\usepackage{subcaption}
\usepackage{booktabs}
\usepackage{array}
\usepackage{tabularx}

% Code listings
\usepackage{xcolor}
\usepackage{listings}

% Links and references
\usepackage[hidelinks]{hyperref}
\usepackage[nameinlink,noabbrev]{cleveref}

% Default epsilon
\let\epsilon\varepsilon

% Common number sets
\newcommand{\R}{\mathbb{R}}
\newcommand{\C}{\mathbb{C}}

% Common quantum/math operators
\newcommand{\id}{\mathbb{I}}
\newcommand{\tr}{\operatorname{tr}}
\newcommand{\Tr}{\operatorname{Tr}}
\newcommand{\diag}{\operatorname{diag}}
\newcommand{\rank}{\operatorname{rank}}

% Code listing defaults
\lstset{
    basicstyle=\ttfamily\small,
    breaklines=true,
    frame=single,
    numbers=left,
    numberstyle=\tiny,
    tabsize=4,
    showstringspaces=false
}

\title{Technical Documentation \--- QIST Project:\\ Cooling
Algorithms for Quantum State Preparation}
\author{Gabriel Maresch}
\date{\today}

\begin{document}

\maketitle
\begin{abstract}
In Ding et al. a cooling protocol is proposed and realized as an averaged quantum channel. The convergence of repeated applications of this channel to a thermal state resp. the ground state is shown, providing end-to-end guarantees. The dynamics of this cooling channel are analyzed by relating them to a Lindbladian approximation. Theorems in Ding show that there exist parameter choices such that (1) the Lindbladian approximately fixes the thermal resp. ground state and (2) the Lindbladian approximates the cooling channel sufficiently well, thus implying that also the fixed point of the cooling channel is close to the desired target state. The actual parameters, however, turn out to be both hard to find, as they have to satisfy a self-consistency condition involving the mixing time, and to suffer from the precondition (2), which effectively forces the coupling to be weak.

In follow-up work of Wang and Ding this restriction is substantially relaxed by deriving state-preservation guarantees directly for the system-bath interaction beyond the weak-coupling Lindblad approximation. In this work we numerically investigate the cooling channel for finite transverse-field Ising models, including parameter regimes beyond the weak-coupling regime of the original analysis. We find evidence that stronger coupling can maintain a fixed point close to the thermal state while requiring far fewer iterations than in the weak-coupling regime. Our numerical results are consistent with the theoretical picture developed in the follow-up work and allow us to study quantitatively how fixed-point accuracy and convergence behavior change as the coupling strength is increased.

\end{abstract}

\section*{Introduction}

Quantum mechanics connects with the tangible physical world by predicting  numerical values of actual measured quantities, i.e. by providing expectation values $\langle \mathcal{O} \rangle$ for the observables of interest. These expectation values only make sense once the state of the system, either pure or mixed is fixed. A well known result from statistical mechanics is, that for systems that are in thermal equilibrium, the states under consideration are precisely the Gibbs states $\rho_{\beta}$ in the case of finite temperature or the ground-state $\ket{\Psi_0}$ in the case of zero-temperature. 

In particular, the time-ordered (resp. retarded) Green's-function of a many body Hamiltonian $H$ at a given inverse temperature $\beta$ is the expectation value for the \emph{correlation} functions
\begin{align*} - i\langle \mathcal{T} c_i(t) c_j^{\dagger}(t')\rangle_{\beta}=
-i \mbox{Tr}\, \rho_{\beta}\, \mathcal{T} c_i(t) c_j^{\dagger}(t')  &\quad {\rm resp.} &- i\langle \{c_i(t) c_j^{\dagger}(t')\}\rangle_{\beta}=-i \mbox{Tr}\, \rho_{\beta}\, \theta(t-t')\{c_i(t), c_j^{\dagger}(t')\}.
\end{align*}

Once the (site dependent) Green's-function $G_{i,j}^T(t-t')$ (resp. $G_{i,j}^R(t-t')$) is known, several important one-particle properties can be effectively computed, e.g.
\begin{description}
    \item[\qquad Spectral function:] $A_{ij}(\omega)=-\frac{1}{\pi} {\rm Im } G_{i,j}^R(\omega)$ 
    \item[\qquad Density of states:] $D(\omega) = -\frac{1}{\pi N}   \mbox{Tr}\, {\rm Im }  G_R(\omega)$
    \item[\qquad Particle density:] $n_i = - i G_{ii}^T(0^-)$
    \item[\qquad Self energy:] $\Sigma(\omega) = G_0(\omega)^{-1}-G(\omega)^{-1}$.
\end{description}

One way to obtain thermal states is not by using some truncation or approximation scheme, but to engineer a specific interaction with a bath, that effectively enforces thermalization at a given temperature. Such a system has the advantage of being easily implementable on a quantum computer or quantum simulator, and allow for running times that often are merely polynomial in the system's parameters. 

The total system's Hamiltonian can be split in three parts:
\begin{align}\label{eq:hamiltonian_parts}
    H_{\rm cooling} = H_0 + H_{\rm int} + H_{\rm bath} 
\end{align}
and a cooling protocol would typically look like
\begin{align}\label{eq:qu_channel}
\rho\otimes\rho_b \longrightarrow U(t)(\rho\otimes\rho_b)\, U^{\dagger}(t)=\Tilde{\rho}  \longrightarrow \Tr_{\mathcal H_{\rm bath}} \Tilde{\rho}  =: \Lambda[\rho]
\end{align}
which corresponds to the following steps:
\begin{enumerate}
    \item prepare both the systems that is to be cooled and the bath: the first in a state $\rho$ and the latter in a thermal state $\rho_b$.
    \item  let the full system evolve for some time $t>0$. Here $U(t) = e^{-i t H}$ for time-independent Hamiltonians resp. $U(t)=\mathcal{T} \exp\left(-i \int_{0}^{ t} dt' H_{\rm cooling} (t') \right)$ for the general case. Also note, that we are working in natural units where $\hbar = 1$.
    \item discard the bath-subsystem.
\end{enumerate}
Note, that the last step will in general produce a density matrix rather than a pure state even if the initial state $\rho$ has been pure. Thus, if we think of repeated application of this scheme, all involved states will be non-pure. The restricted action of \autoref{eq:qu_channel} on the isolated state-space $\mathcal{H}_0$ is usually called a quantum channel ($\Lambda[\,\cdot\,]$) and exhibits richer dynamics than a purely unitary evolution. 

In this project work we report on the implementation of the cooling scheme proposed in \cite{ding2025} which follows along the lines of \autoref{eq:qu_channel} and validate it against a known system Hamiltonian, the TFIM (transverse field Ising model) Hamiltonian and small system sizes $N=4-8$, where exact diagonalization is still possible and the thermal target state 
\begin{align}\label{eq:gibbs_state}
\rho_{\beta}= \frac{e^{-\beta H_0}}{{\rm Tr}\, e^{-\beta H_0}}
\end{align}
may explicitly be calculated.  In this situation $H_0 = H_{\rm TFIM}^{(h,J)}$ is actually a $2^N \times 2^N$ matrix, where $(h,J)$ are the parameters of the model (see \autoref{sec:TFIM}). The bath can be taken as generic two-level system at inverse temperature $\beta$, and its Hamiltonian $H_{\rm bath}$  is just a $2\times 2$ matrix. The interaction Hamiltonian $H_{\rm int}^{(\alpha, A, f, g)}$ is  a block off-diagonal $2^{N+1} \times 2^{N+1}$ matrix that couples the two systems (again $\alpha, A, f, g$ are specific parameters of the coupling and will be explained later).


\section{Cooling as quantum operation}
The two-level system utilized as bath is fully described by its eigenenergies: here they are taken as $\pm \tfrac{\omega}{2}$, where $\omega>0$ is a frequency-parameter. The fiduciary state for the bath subsystem is the two-dimensional Gibbs-state $\rho_{\beta}\propto e^{\tfrac{\omega}{2}\beta Z }$.

The basic building blocks for the coupling terms are of the form
\begin{align}\label{eq:single_cooling_process}
    A_S \otimes \sigma_+ +A_S^{\dagger} \otimes \sigma_- 
\end{align}
each of which describes the physical process of transferring energy from the system to the bath and vice versa. There are a few requirements, both technical and heuristic on how the operator $A_S$ is chosen. Most importantly $A_S$ itself should be treated as parameter and is chosen from pre-specified set of operators $\mathcal{A}$. To de-clutter the algorithms, $\mathcal{A}$ is usually chosen to be invariant under Hermitian conjugation and sign inversion. Furthermore $\|A\|\le 1$, in particular the spectra should be universally bounded. It is essential that $\mathcal{A}$ contains operators that are not diagonal in the energy-eigenbasis and but connects states from different energies, making cooling possible in the first place. In the case that $H_0$ exhibits different symmetry sectors, it is essential that $\mathcal{A}$ mixes those sectors, e.g. if we have a translation-invariant chain of spins, the operators from $\mathcal{A}$ should address the sites individually.  

The term \autoref{eq:single_cooling_process} is multiplied with $\alpha f(t)$, where $\alpha >0$ is fixed and $f(t)$ a time-dependent coupling strength. We require $\int dt \, |f(t)|^2  = 1$, so $f(t)$ resp. its Fourier-transform $\hat{f}(\omega)$ should be thought of as frequency filter and is usually taken as square-root of a Gaussian $\mathcal{N}(0,\sigma^2)$. $\hat{f}(\omega)$ suppresses transitions from the system to the bath and vice versa that correspond to a maximally allowed detuning of Bohr-frequencies from the bath frequency on a scale set by $\Delta \omega \sim \frac{1}{\sigma}$. Thus $H_{\rm int}(t) =  H_{\rm int}^{(\alpha, A, \sigma)}(t)$. 

For the time evolution a duration $T>0$ is fixed and the fiduciary states are evolved from $-T$ to $T$. Note, that the negative times are not essential and could easily be avoided by redefining $f$. However it is beneficial to preserve symmetry of $f$ via $f(-t)=f(t)$ over the time interval $[-T,T]$. The full Hamiltonian, however, is now time dependent which makes the computation of the time-propagation more intricate:
\begin{align}\label{eq:timedependent_H}
U(t_0 \rightarrow t_1) = \mathcal{T}\exp\left(-i \int_{t_0}^{ t_1} d\tau \left(H_0 +  \alpha f(\tau)( A_S \otimes \sigma_+ +A_S^{\dagger} \otimes \sigma_- )+
H_{\rm bath}\right) \right)
\end{align}


While an analog simulation would produce the correct time evolution intrinsically, on a digital quantum simulator we have to program the time evolution explicitly and do so with second order Trotterization for some time $\Delta t$
\begin{align}\label{eq:trotterization}\notag
U(t\rightarrow t+\Delta t)&\approx \exp\left(-i\int_{t+\tfrac 12 \Delta t}^{t+\Delta t}d\tau H_{\rm int}(\tau)\right)e^{-i(H_0+H_{\rm bath})\Delta t} \exp\left(-i\int_{t}^{t+\tfrac 12 \Delta t}d\tau H_{\rm int}(\tau)\right)\\
\end{align}
and then approximate the integrals by using the value at the midpoint as constant approximation in order to preserve symmetry:
\begin{align}\label{eq:trotterization_with_freezing}
U(t\rightarrow t+\Delta t)&\approx e^{-i H_{\rm int}(t+\tfrac 12 \Delta t)\tfrac 12 \Delta t}e^{-i(H_0+H_{\rm bath})\Delta t} e^{-i H_{\rm int}(t+\tfrac 12 \Delta t)\tfrac 12 \Delta t}.
\end{align}
In the cooling algorithm itself, the interval $[-T,T]$ is divided into $M$ subintervals of length $\tau$ and each subinterval is evolved according to \autoref{eq:trotterization_with_freezing}. Note, that freezing $H_{\rm int}$ at the subinterval-midpoint is just a convenient numerical choice, but not necessarily optimal. Using different evaluation points for the quadrature of $f$, or even exact evaluation of the $f$-integration, would also be a second order approximation, but with potentially reduced constant of the error term. However, we will not pursue this particular direction further.

\section{Transverse field Ising model}\label{sec:TFIM}
The benchmark system which is used for validation is given by the Hamiltonian
\begin{align*}
    H(J,h) = -J\sum_{<i,j>} \sigma^{(z)}_{i}\sigma^{(z)}_{j} - h\sum_{i=0}^{N-1} \sigma^{(x)}_i 
\end{align*} 
and models a 1D-chain of spins in $z$-direction in a homogeneous magnetic field $h$ that is aligned transversally with respect to the orientation of the spins. For simplicity we only consider periodic boundary conditions and an even number of sites. The relevant parameter of this system is $J/h$, the ratio of the energy scale of the spin-spin interaction to the potential energy of the spins in the magnetic field. To keep the energy-scales of $H(J,h)$ comparable while sweeping $J$ and $h$, we usually will impose the normalization
\begin{align*}
    J \to J/\sqrt{J^2+h^2}\\
    h \to h/\sqrt{J^2+h^2}
\end{align*} which leaves the underlying physics unchanged. For $J\ll h$ a configuration with all spins aligned in $\ket{\uparrow}_x$ is favored (paramagnetic phase) and the spectrum is gapped with $\Delta \approx 2|h|$. 

For $J\gg h$ the ground-state is twofold degenerate where for $J>0$ all spins are either aligned in $\ket{\uparrow}_z$ direction or in $\ket{\uparrow}_z$ direction (ferromagnetic phase) resp. aligned alternately for $J<0$. The gap to the first excited state is $\Delta \approx 2 |J|$

The gaps remain finite as long as the order parameter $J/h$ stays away from the quantum critical region $J\approx h$. Here the gap closes. 

The spectrum of $H(J,h)$ can be computed exactly by mapping the TFIM to a free fermionic system
\begin{align*}
    H = \sum_k \varepsilon_k (J,h)\left( \hat{n}_k - \frac12\right), \quad
    \varepsilon_k(J,h) = 2\sqrt{(h-J\cos k)^2 + J^2\sin^2 k}.
\end{align*}
In particular for $J=h$ one gets
\begin{align*}
    H = \sum_k 4J \big| \sin \frac{k}{2}\big| \left( \hat{n}_k - \frac12\right)
\end{align*} which turns out to exhibit conformal symmetry in the limit for $k\approx 0$ and $N\to \infty$, however we will not make use of this fact in the present work. The equivalence of TFIM to a free fermionic Hamiltonian on the other hand allows for a-priori estimates on the integer mixing time of the cooling algorithm, that would be much harder to come by for other models.

\section{Parameter dependance}
Let us quickly review the parameters we encountered so far:

\begin{description}
    \item[coupling strength $\alpha$] determines the magnitude of the interaction part in the cooling Hamiltonian, i.e. the maximal rate at which cooling processes mediated by the bath can take place. To allow for the description of the cooling process via perturbative Lindblad dynamics $\alpha$ is treated as small parameter in \cite. However, larger coupling allows for higher energy transfer from the system to the bath and therefore increases cooling. In \cite the validity of the cooling scheme is established also for $\alpha = O(1)$. The numerical experiments of this work underscore these findings.  
    \item[frequency cutoff $\sigma$], where $\sigma$ denotes standard deviation of a Gaussian function in frequency space. Tunes scale at which transitions with Bohr-frequencies bigger than $\sigma$ are suppressed. A safe choice is 
    $\|H_0\|$ resp. the energy bandwidth of the system that is to be cooled.
    
\end{description}

\section{Implementation as quantum circuit}
In this project the pennylange python package was used to implement concrete realizations of the quantum channel. Here 'concrete' means, that the total cooling Hamiltonian $H$ is built using a specific jump-operator $A\in \mathcal{A}$ and a specific Bohr frequency $\omega\in [0, \omega_0]$ and evolved from time $t=-T$ to time $t=T$. The second order Trotterization from eq (4) is directly implemented into the circuit, meaning
once the total evolution time $2T$ and the substep time $\tau$ is chosen, the quantum circuit is built according to this specifications. It consists of $M+1$ layers, each effecting a time increment of $\tau$. In \autoref{}
a circuit which implements four cooling steps each of which is uses individual second order Trotter steps of length $\tau = T/2$. Here $W(\tau)$ denotes the time-evolution of $H_{\rm int}$, whereas the time evolution of $H_{\rm bath}$ is implemented directly as $Z$-rotation and the time evolution of $H_0$ is obtained by second order Trotterization with the same step size $\tau$.

After each cooling step, the auxilary qubit is measured and prepared in the same Gibbs-state of the two-level system. Pennylane offers two approaches on how to implement this reset:

\begin{itemize}
    \item we choose the inital state of the auxilary qubit probabilistically according to the statistics of the Gibbs state, i.e. 
    \begin{align}
        p_0=p(\ket{0}) = \frac{e^{-\beta \tfrac{\omega}{2}}}{2\cosh{\beta \tfrac{\omega}{2}}} \quad\mbox{resp.}\quad p_1=p(\ket{1}) = \frac{e^{\beta \tfrac{\omega}{2}}}{2\cosh{\beta \tfrac{\omega}{2}}}.
    \end{align}
    \item we run the simulation in mixed mode, meaning states are not assumed to be pure but mixed instead. In this case we can implement the constant 2-dimensional quantum channel 
    \begin{align*}
        \rho \mapsto \sum_{i,j} p_i \underbrace{\ket{i}\bra{j}}_{K_{ij}} \rho \underbrace{\ket{j}\bra{i}}_{K_{ij}^{\dagger}} = \sum_{i} p_i \ket{i}\bra{i}\, \rm{Tr}\rho = \rho_{\beta \tfrac{\omega}{2}},
    \end{align*} where we have used the channel representation with Kraus-operators, which can be implemented directly in pennylane.
\end{itemize}

Nota bene: Simulation of the quantum circuit makes it necessary to decide in advance how many iteration steps shall be undertaken. In a physical quantum computer the gate sequence which constitutes one cooling step, would instead be applied iteratively (e.g. a fixed sequence of laser pulses) in a trapped ion or neutral atom setup.

Furthermore, since for simulation all iteration steps have to be pre-determined, meaning a specific sequence $(A_n,\omega_n)\in \mathcal{A}\times [0, \omega_{\rm max}]$ has to be chosen in advance. In the present implementation this is done with a random seed, in order to allow for reproducability.


\section{Simulation using superoperator representation}
A quantum channel is a map that maps density matrices to density matrices, i.e. $\Lambda: \mathbb{C}^{2^N\times 2^N} \to \mathbb{C}^{2^N\times 2^N}$. If we view $\Lambda$ as a transformation of matrices, it is common to use the name superoperator instead. $\Lambda$ can then be treated as a $4^N\times 4^N$ matrix. The identification $\mathbb{C}^{2^N\times 2^N}\cong \mathbb{C}^{4^N}$ is called ,,vectorization''.

\subsection{Choi representation}

\subsection{Kraus representation}
Every quantum channel has also - at least one - Kraus representation, i.e. can be written in the form
\begin{align*}
    \Lambda[\rho] = \sum_i K_i \rho K_i^{\dagger}.
\end{align*}

Since the systems's state space $\mathcal{H}=\mathcal{H}_0\otimes \mathcal{H}_b$, where the bath's state space is just two-dimensional, this allows for some simplification: Let us denote by $U=U(t_m,t_{m+1}$ the time evolution of composite system including bath for a particular time interval. We can write
\begin{align*}
U =
\begin{pmatrix}
U_{00} & U_{10} \\
U_{01} & U_{11},
\end{pmatrix}
\end{align*}
where $U_{ij} = \bra{i}_{b} U  \ket{j}_{b}$ is an operator $\mathcal{H}_0 \to \mathcal{H}_0$, i.e. a $2^N \times 2^N$ matrix. Applying $U\rho U^{\dagger}$ to a density matrix of the form $\rho_0 \otimes \sum_{i} p_i  \ket{i}_{b}\bra{i}_{b}$ we get
\begin{align*}
     \Tr_{\mathcal H_{\rm bath}} U\rho U^{\dagger} = \sum_{ij} p_j  U_{ij}\rho_0 U_{ij}^{\dagger}.
\end{align*} This involves only matrix-products of $2^N \times 2^N$ matrices, which is advantageous to multiplying $2^{N+1} \times 2^{N+1}$ matrices.

\subsection{Averaged channel}
The cooling algorithm proposed by Ding et.al. is probabilistic in the sense, that end-to-end guarantees apply to the \emph{average} channel, where an expectation value is taken both with respect to the jump operators and the Bohr frequencies:
\begin{align*}
    \overline{\Lambda}\,[\rho]:= \mathbb{E}_{A, \omega} \Lambda(A,\omega)\,[\rho]. 
\end{align*}

In choosing the underlying probability distributions on $\mathcal{A}$ resp. $[0,\omega_{\max}$ we follow Ding et. al. and choose uniform distributions on both sets. At least for $\mathcal{A}$ this seems to be the canonical choice - at least if the jump operators correspond to physical sites and the Hamiltonian is translation invariant. 

In this work we focus on the set $\mathcal{A}=\{\pm \sigma_i^{(x)}, \pm \sigma_i^{(z)}: i = 1,\ldots N\}$ and the uniform probability assigning each operator $\frac{1}{4N}$.
As simple choice for the set of Bohr frequencies $[0,\omega_{\rm}]$ we choose uniform distribution as well. Here however there is a-priori no need to exclude other distributions and it could turn out that e.g. distributions 
that are skewed towards higher frequencies have an advantage. In this work however we don't pursue this particular degree of freedom.

To realize the averaged channel $\overline{\Lambda}$ numerically we approximate
\begin{align*}
    \int_0^{\omega_{\rm max}} d\omega\, \Lambda(A,\omega) = \frac{1}{k}\sum_{i=1}^k \Lambda(A,\omega_i) 
\end{align*} by some quadrature formula, e.g. by the midpoint-formula where $\omega_i = \frac{2i+1}{k}\omega_{\rm max}$. The average channel is thus approximated by the object
\begin{align*}
\frac{1}{4N\cdot k \omega_{\rm max} }\sum_{A\in\mathcal A} \sum_{i=1}^k \Lambda(A, \tfrac{2i+1}{k}\omega_{\rm max})
\end{align*}


up to order $O(\tfrac{1}{k^2})$. Thus the number of individual contributions for a Hamiltonian on $N$ sites scales with $4N\cdot k$. To keep the computational demand in check we let $k = 10$, so for computating the average channel for $N=4,6,8$ we have to run trajectories of length $160, 240$ and $320$ each. 
\pagebreak

\section{Numerical Experiments}
\subsection{h-sweep}


\begin{figure}[htbp]
    \centering
    \begin{subfigure}[t]{0.49\textwidth}
        \centering
        \includegraphics[width=\linewidth]{docu/iterations_normalized.png}
        \caption{First subplot.}
        \label{fig:subplot1}
    \end{subfigure}%
    \hfill
    \begin{subfigure}[t]{0.49\textwidth}
        \centering
        \includegraphics[width=\linewidth]{docu/iterations_log_normalized.png}
        \caption{Second subplot.}
        \label{fig:subplot2}
    \end{subfigure}

    \caption{Overall caption.}
    \label{fig:subplots}
\end{figure}

\begin{figure}[htbp]
    \centering
    \includegraphics[width=0.8\textwidth]{docu/gap_normalized.png}
    \caption{Description of the image.}
    \label{fig:image-label}
\end{figure}

\begin{figure}[htbp]
    \centering
    \includegraphics[width=0.8\textwidth]{docu/distance_to_gibbs_normalized.png}
    \caption{Description of the image.}
    \label{fig:image-label}
\end{figure}

\subsection{$\sigma$-sweep}
\subsection{$\beta$-sweep}



\section{Results}



\begin{thebibliography}{9}

\bibitem{ding2025}
Z.~Ding, Y.~Zhan, J.~Preskill, and L.~Lin,
\textit{End-to-End Efficient Quantum Thermal and Ground State Preparation Made Simple},
arXiv:2508.05703 [quant-ph] (2025).

\bibitem{wang2025}
K.~Wang and Z.~Ding,
\textit{Beyond Lindblad Dynamics: Rigorous Guarantees for Thermal and Ground State Preservation under System Bath Interactions},
arXiv:2512.03457 [quant-ph] (2025).

\end{thebibliography}

\end{document}
